\chapter{METHODOLOGY}

\section{Introduction}
This chapter describes the research design, mathematical formulations, and 
experimental setup utilized to implement and evaluate the proposed 
workload-aware dynamic array growth algorithm. To systematically develop an 
adaptive resizing framework, we constructed a custom Python-based benchmarking 
simulation environment that models different data ingestion patterns. 

The research is organized into four progressive phases: (i) mathematical 
formulation of the adaptive growth functions, (ii) simulator design and workload 
generation, (iii) parameter optimization using Grid Search and Bayesian 
Optimization, and (iv) performance evaluation against static industry baselines. 
By measuring memory waste, element copying overhead, and execution times, we 
quantify the efficiency of the proposed algorithms under identical, controlled 
workloads of up to $10^7$ sequential insertions.

\section{Materials}
The implementation and empirical evaluations were conducted in a standardized 
software and hardware environment to ensure reproducibility and consistency 
across all test trials.

\subsection{Software Environment}
The simulation engine was developed in \textbf{Python 3.11.4 (64-bit)}. Python 
was selected due to its rapid prototyping capabilities, high-level memory 
management abstractions, and the availability of scientific computing libraries. 
The core packages used in this research include:
\begin{itemize}
    \item \textbf{Scikit-Optimize (\texttt{skopt} v0.9.0):} Used for implementing 
    the Gaussian Process (GP) Bayesian Optimization framework. Specifically, the 
    \texttt{gp\_minimize} function was utilized to find the optimal hyper-parameters 
    for the growth algorithms.
    \item \textbf{NumPy (\texttt{numpy} v1.24.3):} Employed for fast numerical array 
    operations and generating synthetic workloads (such as sinusoidal and random 
    burst velocities).
    \item \textbf{Matplotlib (\texttt{matplotlib} v3.7.1):} Used to visualize 
    parameter sensitivity curves, convergence traces, and relative performance 
    histograms.
\end{itemize}

\subsection{Hardware Environment}
All benchmarks, grid sweeps, and Bayesian optimization loops were executed on 
a single, dedicated workstation with the following hardware specifications to 
prevent CPU throttling and external measurement noise:
\begin{itemize}
    \item \textbf{CPU:} Intel Core i5-11400F @ 2.60GHz (6 Cores, 12 Threads, 12MB Cache)
    \item \textbf{RAM:} 16GB DDR4 @ 3200MHz
    \item \textbf{Operating System:} Windows 10 Pro (64-bit)
    \item \textbf{Storage:} 512GB NVMe M.2 SSD
\end{itemize}
To ensure time-measurement accuracy, all background processes were disabled during 
benchmarking, and execution times were captured using the high-resolution 
monotonic clock interface (\texttt{time.perf\_counter()}) which provides 
nanosecond-level precision.

\section{Research Design}

\subsection{Mathematical Formulations of Growth Families}
To systematically analyze adaptive resizing, we formulated four distinct 
growth families, each introducing progressive degrees of workload awareness. 
To prevent mathematical collapse (such as growth factors falling below 1.0 or 
climbing to values that trigger memory overflow), all growth factors are 
strictly clamped using a bounding function:
\begin{equation}
    g_{\text{clamped}}(k, v) = \min\left(g_{\max}, \max\left(g_{\min}, g(k, v)\right)\right)
\end{equation}
where $g_{\min} = 1.1$ and $g_{\max} = 2.0$ represent the lower and upper guard 
rails, respectively.

\subsubsection{Formula Family A: Logarithmic Decay}
Formula A models the growth factor as a decaying logarithmic function of the 
array's capacity. The core theory is that as the array size increases, the 
absolute amount of memory wasted grows exponentially under fixed geometric 
expansion. To counter this, the growth factor decays logarithmically.
\begin{itemize}
    \item \textbf{Phase 1 (Static Logarithmic Growth):}
    \begin{equation}
        g_A(C_{\text{old}}) = 1 + \frac{\alpha}{\ln(C_{\text{old}} + \beta)}
    \end{equation}
    where $C_{\text{old}}$ is the current array capacity, $\alpha$ is a scaling 
    coefficient, and $\beta$ is a shifting factor that prevents division by zero.
    \item \textbf{Phase 2 (Workload-Aware Logarithmic Growth):}
    \begin{equation}
        g_A(C_{\text{old}}, v) = 1 + \frac{\alpha + \lambda \ln(1 + v)}{\ln(C_{\text{old}} + \beta)}
    \end{equation}
    where $v$ is the insertion velocity, and $\lambda$ is a weight scaling the 
    influence of velocity on array expansion.
\end{itemize}

\subsubsection{Formula Family B: Harmonic Decay}
Formula B models the growth factor as a function of the resize iteration history, 
denoted as $k$. As the array undergoes more resizes, the growth factor decays 
following the harmonic series:
\begin{itemize}
    \item \textbf{Phase 1 (Static Harmonic Growth):}
    \begin{equation}
        g_B(k) = 1 + \frac{c_0}{k}
    \end{equation}
    where $k$ is the current resize count ($k = 1, 2, 3...$) and $c_0$ is the 
    base growth multiplier.
    \item \textbf{Phase 2 (Workload-Aware Harmonic Growth):}
    \begin{equation}
        g_B(k, v) = 1 + \frac{c_0 + \lambda \ln(1 + v)}{k}
    \end{equation}
    where $c_0$ controls the baseline expansion rate, and $\lambda \ln(1+v)$ 
    adjusts the expansion dynamically based on the log-velocity of data 
    ingestion.
\end{itemize}

\subsubsection{Formula Family C: Manual Hybrid Framework}
Formula C is a manual linear combination of the logarithmic (Formula A) and 
harmonic (Formula B) designs. It attempts to merge the capacity-dependent decay 
of Formula A with the step-wise decay of Formula B using static weights:
\begin{equation}
    g_C(k, v, C_{\text{old}}) = w_A \cdot g_A(C_{\text{old}}, v) + (1 - w_A) \cdot g_B(k, v)
\end{equation}
where $w_A \in [0.0, 1.0]$ is a manually assigned weight defining the balance 
between the two models.

\subsubsection{Formula Family D: GP-Optimized Hybrid Framework}
Formula D represents the generalized workload-aware harmonic expansion model. 
Unlike Formula C, which manually blends two distinct structures, Formula D 
focuses strictly on the harmonic formulation (which proved superior in baseline 
trials) but subjects all hyper-parameters to automated Gaussian Process (GP) 
Bayesian Optimization:
\begin{equation}
    g_D(k, v) = \min\left(g_{\max}, \max\left(g_{\min}, 1 + \frac{c_0 + \lambda \ln(1 + v)}{k}\right)\right)
\end{equation}
Here, the parameters $c_0$ and $\lambda$ are optimized globally to minimize the 
unified cost score $J$ across mixed workloads.

\section{Data Collection}

\subsection{Workload Ingestion Profiles}
To test the adaptability of the algorithms, we designed seven distinct workload 
profiles that simulate realistic ingestion patterns:
\begin{enumerate}
    \item \textbf{Uniform Workload:} Elements are inserted at a constant, steady 
    rate. This represents standard static streaming. Mathematically, the 
    number of elements $N(t)$ at time step $t$ is expressed as:
    \begin{equation}
        N(t) = \text{round}(C_{\text{uniform}} \cdot t)
    \end{equation}
    where $C_{\text{uniform}}$ is the constant insertion rate.
    \item \textbf{Bursty Workload (Spike Ingestion):} Alternating cycles of 
    extended inactivity and extremely rapid, high-volume element appends. The 
    insertion rate oscillates between zero and a peak value:
    \begin{equation}
        N'(t) = \begin{cases} 
            V_{\text{burst}} & \text{if } t \pmod{T_{\text{period}}} < T_{\text{active}} \\
            0 & \text{otherwise}
        \end{cases}
    \end{equation}
    where $T_{\text{period}}$ is the workload period and $T_{\text{active}}$ is the 
    duration of the active burst phase.
    \item \textbf{Mixed Workload:} A combination of uniform phases and burst phases 
    to replicate complex runtime environments.
    \item \textbf{Increasing Workload:} The rate of element insertion increases 
    exponentially over time. The cumulative element count $N(t)$ is defined as:
    \begin{equation}
        N(t) = N_0 \cdot e^{r \cdot t}
    \end{equation}
    where $N_0$ is the initial insert count and $r$ is the growth rate constant.
    \item \textbf{Decreasing Workload:} A high initial insertion rate that 
    tapers off to a slow trickle. The cumulative element count $N(t)$ is defined as:
    \begin{equation}
        N(t) = N_{\max} \cdot \left(1 - e^{-r \cdot t}\right)
    \end{equation}
    where $N_{\max}$ is the asymptotic upper limit of the array size.
    \item \textbf{Sinusoidal Workload:} Ingestion rate oscillates periodically, 
    representing cyclical traffic (e.g., daytime vs. nighttime loads). The instantaneous 
    rate is:
    \begin{equation}
        N'(t) = A \cdot \sin\left(\frac{2\pi}{T} \cdot t\right) + B
    \end{equation}
    where $A$ is the amplitude, $T$ is the cycle period, and $B$ is the baseline rate.
    \item \textbf{Step Workload:} Abrupt shifts in insertion rates, changing 
    from one steady state to another.
\end{enumerate}

\subsection{Insertion Velocity Metric}
The insertion velocity $v$ is calculated dynamically during each insertion 
cycle. When size reaches capacity, the instantaneous velocity is calculated 
before triggering reallocation:
\begin{equation}
    v = \frac{\Delta N}{\Delta t + \epsilon}
\end{equation}
where $\Delta N$ is the number of elements inserted since the last resize event 
(which is always equal to $C_{\text{old}} - C_{\text{prev\_old}}$), and $\Delta t$ 
is the elapsed wall-clock time between the previous resize and the current 
resize. The small regularization constant $\epsilon = 10^{-9}$ is added to 
prevent division-by-zero errors in ultra-fast environments.

\section{Data Analysis: The Unified J-Score}
To evaluate and compare configurations, we developed a multi-objective cost 
function, the $J$-Score. Because minimizing memory waste often increases 
element copies, and vice-versa, a single scalar score is required to find the 
mathematical Pareto-optimum. The unified objective function is defined as:
\begin{equation}
    J = w_{\text{waste}} \cdot \overline{W} + w_{\text{copies}} \cdot \overline{C_{\text{copy}}} + w_{\text{time}} \cdot \overline{T}
\end{equation}
Specifically, the weights were set to prioritize memory efficiency and 
computational copying overhead equally, with execution time serving as a 
tertiary constraint:
\begin{equation}
    J = 0.4 \cdot \overline{W} + 0.4 \cdot \overline{C_{\text{copy}}} + 0.2 \cdot \overline{T}
\end{equation}
Here, the components are normalized ratios:
\begin{itemize}
    \item $\overline{W}$ is the normalized memory waste ratio, calculated as 
    the wasted capacity at the end of the workload divided by the final capacity.
    \item $\overline{C_{\text{copy}}}$ is the normalized cumulative element copies 
    relative to a worst-case baseline policy.
    \item $\overline{T}$ is the normalized total execution time relative to the 
    fixed $2.0\times$ baseline.
\end{itemize}
A lower $J$-Score indicates a more efficient growth strategy.

\section{Implementation}

\subsection{Simulator Engine Architecture}
The benchmarking simulator mimics a contiguous dynamic array. It maintains a 
virtual contiguous address space, an active size, and an active capacity. When 
elements are appended, the simulator increases the size. If the size exceeds 
capacity, it queries the growth engine, evaluates the current velocity $v$ and 
resize counter $k$, computes the new capacity $C_{\text{new}} = \lfloor C_{\text{old}} \times g(k,v) \rfloor$, 
allocates new virtual memory, copies the old elements, and frees the previous 
allocation.

\begin{figure}[H]
\centering
\begin{tikzpicture}[
    block/.style={draw, fill=blue!10, rectangle, minimum width=2.5cm, minimum height=1.0cm, align=center, font=\small},
    decision/.style={draw, fill=yellow!15, diamond, minimum width=2.5cm, minimum height=2.0cm, align=center, font=\small},
    io/.style={draw, fill=green!10, trapezium, trapezium left angle=70, trapezium right angle=110, minimum width=2cm, minimum height=1cm, align=center, font=\small},
    arrow/.style={-Stealth, thick}
]
    % Nodes
    \node[io] (start) {Initialize Array \\ $C_0 = 10, k = 0, \text{Size} = 0$};
    \node[block, below=0.8cm of start] (insert) {Insert Element \\ from Workload Pattern};
    \node[decision, below=0.8cm of insert] (check) {Is Array \\ Full? \\ $\text{Size} = C_{\text{old}}$};
    \node[block, right=1.5cm of check] (no_resize) {Increment Size \\ Continue Loop};
    \node[block, below=0.8cm of check] (calc_v) {Calculate Velocity \\ $v = \frac{\Delta N}{\Delta t}$};
    \node[block, below=0.8cm of calc_v] (calc_g) {Compute Clamped $g(k,v)$ \\ Formula A/B/C/D};
    \node[block, below=0.8cm of calc_g] (resize) {Allocate $C_{\text{new}} = \lfloor C_{\text{old}} g \rfloor$ \\ Copy Elements, Free Old \\ Increment $k$};
    \node[decision, below=0.8cm of resize] (done) {Workload \\ Completed?};
    \node[io, below=0.8cm of done] (end) {Return Logged \\ Metrics (J-Score)};

    % Arrows
    \draw[arrow] (start) -- (insert);
    \draw[arrow] (insert) -- (check);
    \draw[arrow] (check) -- node[above] {No} (no_resize);
    \draw[arrow] (no_resize) |- (insert);
    \draw[arrow] (check) -- node[left] {Yes} (calc_v);
    \draw[arrow] (calc_v) -- (calc_g);
    \draw[arrow] (calc_g) -- (resize);
    \draw[arrow] (resize) -- (done);
    \draw[arrow] (done) -- node[left] {Yes} (end);
    \draw[arrow] (done.west) -- node[below] {No} ++(-1.5,0) |- (insert.west);
\end{tikzpicture}
\caption{System Architecture Flowchart of the Resizing Simulator Engine}
\label{fig:simulator_flowchart}
\end{figure}

\subsection{Python Simulator Implementation Listing}
To provide absolute clarity on the simulator's computational layout, Listing 
3.1 presents the core implementation of the \texttt{AdaptiveDynamicArray} engine, 
capturing the velocity calculations and metric tracking during execution.

\begin{lstlisting}[language=Python, caption={Core Implementation of the Adaptive Dynamic Array Simulator}, label={lst:simulator_code}]
import time
import math

class AdaptiveDynamicArray:
    def __init__(self, growth_formula_fn, initial_capacity=10):
        self.capacity = initial_capacity
        self.size = 0
        self.resize_count = 0
        self.elements_copied = 0
        
        # Velocity tracking states
        self.last_resize_time = time.perf_counter()
        self.last_resize_size = 0
        
        # Growth factor function hook
        self.growth_formula = growth_formula_fn
        
    def append(self, element):
        if self.size >= self.capacity:
            self._resize()
        self.size += 1

    def _resize(self):
        current_time = time.perf_counter()
        delta_t = current_time - self.last_resize_time
        delta_n = self.size - self.last_resize_size
        
        # Compute velocity (v = delta_n / delta_t)
        velocity = delta_n / (delta_t + 1e-9)
        self.resize_count += 1
        
        # Query growth factor
        g = self.growth_formula(self.resize_count, velocity, self.capacity)
        
        # Clamping to guard rails (1.1x - 2.0x)
        g = min(2.0, max(1.1, g))
        
        # Compute new capacity
        new_capacity = math.floor(self.capacity * g)
        
        # Track metric outputs
        self.elements_copied += self.size
        self.capacity = new_capacity
        
        # Reset velocity trackers
        self.last_resize_time = time.perf_counter()
        self.last_resize_size = self.size
\end{lstlisting}

\section{Evaluation}

\subsection{Grid Search Sweep}
To map the parameter space and establish initial performance baselines, a grid 
search sweep was conducted for Formulas A, B, and C.
\begin{itemize}
    \item \textbf{Formula A Sweep:} $\alpha$ was swept from $0.5$ to $5.0$ with a 
    step of $0.5$, and $\beta$ was swept from $2$ to $20$ with a step of $2$.
    \item \textbf{Formula B Sweep:} $c_0$ was swept from $0.5$ to $4.0$ with a step 
    of $0.5$, and $\lambda$ was swept from $0.0$ to $5.0$ with a step of $0.5$.
    \item \textbf{Formula C Weight Sweep:} The weight $w_A$ was swept from $0.0$ to 
    $1.0$ with a step of $0.1$.
\end{itemize}
Each configuration was executed across 10 independent trials for all 7 workloads, 
generating a total of 15,400 benchmark runs.

\subsection{Gaussian Process Bayesian Optimization Setup}
To locate the global Pareto-optimal parameters for Formula D, we utilized 
Bayesian Optimization using Gaussian Process regression as the surrogate model. 
Bayesian Optimization is suited for optimizing expensive, black-box functions 
by building a probabilistic model of the objective function and using an 
acquisition function to determine where to sample next.

\begin{figure}[H]
\centering
\begin{tikzpicture}[
    node distance=1.5cm and 1.5cm,
    block/.style={draw, fill=blue!10, rectangle, minimum width=2.5cm, minimum height=1.0cm, align=center, font=\small},
    arrow/.style={-Stealth, thick}
]
    % Nodes
    \node[block] (param) {Select Parameters \\ $(c_0, \lambda)$};
    \node[block, right=of param] (sim) {Run Simulator \\ mixed workload};
    \node[block, right=of sim] (jscore) {Evaluate $J$-Score \\ Cost Output};
    \node[block, below=of jscore] (gp) {Update Gaussian Process \\ Surrogate Model};
    \node[block, left=of gp] (acq) {Evaluate Acquisition \\ Function (EI)};
    
    % Connections
    \draw[arrow] (param) -- (sim);
    \draw[arrow] (sim) -- (jscore);
    \draw[arrow] (jscore) -- (gp);
    \draw[arrow] (gp) -- (acq);
    \draw[arrow] (acq) |- ($(param.south) + (0,-0.5)$) -| (param);
\end{tikzpicture}
\caption{The Gaussian Process Bayesian Optimization Feedback Loop}
\label{fig:bo_feedback_loop}
\end{figure}

The optimization space was bounded as follows:
\begin{equation}
    c_0 \in [0.1, 5.0], \quad \lambda \in [0.0, 10.0]
\end{equation}
The objective function to minimize was the composite $J$-Score evaluated under 
the $10^7$ element mixed workload. The acquisition function utilized was the 
\textit{Expected Improvement} (EI):
\begin{equation}
    \text{EI}(x) = \mathbb{E} \left[ \max(0, J(x^+) - J(x)) \right]
\end{equation}
where $x^+$ is the location of the current minimum $J$-Score. 

The Gaussian Process regression modeling utilized a Matern $5/2$ covariance 
kernel function, which is defined as:
\begin{equation}
    K(d) = \sigma^2 \left( 1 + \frac{\sqrt{5}d}{\ell} + \frac{5d^2}{3\ell^2} \right) \exp\left(-\frac{\sqrt{5}d}{\ell}\right)
\end{equation}
where $d = ||x - x'||$ represents the Euclidean distance in hyper-parameter space, 
$\sigma^2$ is the output variance scaling factor, and $\ell$ is the characteristic 
length-scale parameter. The optimization loop was run for $50$ iterations: the first 
$10$ iterations used random initialization to map the surrogate space, followed by 
$40$ sequential guide iterations.

\section{Summary for the Chapter}
This chapter detailed the methodology used to model, implement, and evaluate 
the proposed workload-aware dynamic array growth algorithms. Section 3.2 
presented the software (Python 3.11, NumPy, Scikit-Optimize) and hardware 
specifications of the benchmarking system. Section 3.3 detailed the mathematical 
formulations of the Logarithmic (A), Harmonic (B), Manual Hybrid (C), and 
GP-Optimized Hybrid (D) frameworks, and explained the clamping mechanisms. Section 
3.4 described the 7 synthetic workload profiles, including their exact mathematical 
models, and detailed the velocity metric calculation. Section 3.5 formulated the 
unified multi-objective $J$-Score, weighing memory waste, copying overhead, and 
execution time. Section 3.6 illustrated the simulator's architectural flowchart 
using TikZ and presented the core Python code listing for the simulator engine. 
Section 3.7 detailed the Grid Search and Gaussian Process Bayesian Optimization 
procedures, including the Matern covariance kernel and acquisition functions, 
visualized via a feedback loop diagram. The empirical results obtained using this 
methodology are presented and analyzed in the next chapter.
